\section{Introduction}
\label{sec:introduction}

We prove the Bochner--Riesz conjecture in three dimensions in its
strict-order formulation. The main estimate is boundedness on $L^3$
for every positive order; interpolation and duality then give the full
range of exponents.

For $f\in\mathcal S(\RR^3)$, use the Fourier transform
\[
 \widehat f(\xi)=\int_{\RR^3}f(x)e^{-2\pi i x\cdot\xi}\,dx,
\]
where $\mathcal S$ denotes the Schwartz space. For $\delta>0$, define
\begin{equation}
 T_\delta f(x)
 =\int_{\RR^3}(1-\abs{\xi}^2)_+^\delta
       \widehat f(\xi)e^{2\pi i x\cdot\xi}\,d\xi .
 \label{eq:main-multiplier}
\end{equation}

\begin{theorem}
\label{thm:bochner-riesz}
For every $\delta>0$, there is a finite constant $C_\delta$ such that
\begin{equation}
 \norm{T_\delta f}_{L^3(\RR^3)}
 \le C_\delta\norm f_{L^3(\RR^3)}
 \qquad(f\in\mathcal S(\RR^3)).
 \label{eq:main-bound}
\end{equation}
Consequently $T_\delta$ has a unique bounded extension to $L^3(\RR^3)$.
\end{theorem}

Theorem~\ref{thm:bochner-riesz} gives
\begin{equation}
 \norm{T_\delta f}_{L^p(\RR^3)}
 \lesssim_{p,\delta}\norm f_{L^p(\RR^3)},\qquad
 \delta>\max\left\{3\left|\frac1p-\frac12\right|-\frac12,0\right\},
 \quad 1\le p\le\infty .
 \label{eq:intro-full-range}
\end{equation}
Corollary~\ref{cor:full-bochner-riesz} supplies the interpolation
argument, including the bounds for complex orders. The order is fixed
before choosing any constants. Dilation gives the same bounds for
multipliers $(1-|\xi|^2/R^2)_+^\delta$ at every $R>0$; for finite $p$
in \eqref{eq:intro-full-range}, these means converge to $f$ in $L^p$
as $R\to\infty$. This recovers Fourier inversion by arbitrarily mild
spherical smoothing at the central exponents $3/2\le p\le3$.

\subsection{The problem and its development}

The problem belongs to the theory of summability of Fourier expansions.
Riesz's work on one-dimensional summation methods
\citep{Riesz1923} preceded Bochner's spherical means for multiple
Fourier series \citep{Bochner1935,Bochner1936}. The multiplier
$(1-|\xi|^2)_+^\delta$ softens the boundary of the frequency ball;
the question is how much smoothing is needed to preserve an $L^p$
norm. Radial inversion and radial summability were studied by Herz
and Welland \citep{Herz1954,Welland1975}. For general functions the
geometry of frequency directions is essential. Using a Kakeya
construction, Fefferman proved that in dimension at least two the
order-zero ball multiplier is unbounded on $L^p$ for $p\ne2$
\citep{Fefferman1971}, while Carleson and Sj\"olin established the
conjectured positive-order range in the plane \citep{CarlesonSjolin1972}.

Geometric decompositions connect this problem with the arrangement of
long thin tubes and with Fourier restriction. C\'ordoba's work on
radial multipliers and Kakeya maximal functions developed this
connection through rectangle decompositions and almost orthogonality
\citep{Cordoba1975,Cordoba1977}. Bourgain combined geometric
maximal-function estimates with improved restriction and multiplier
bounds in three dimensions \citep{Bourgain1991RestrictionMultiplier}.
The bilinear restriction theorems of Wolff for the cone and Tao for
the paraboloid and elliptic surfaces
\citep{Wolff2001,Tao2003} led, through Lee's multiplier reduction,
to sharp-order Bochner--Riesz estimates for
$\max\{p,p'\}\ge10/3$ \citep{Lee2004}; here $p'$ is the
H\"older conjugate of $p$. The threshold exponent follows from
Lee's printed strict range by interpolation with $L^2$, keeping a
positive margin in the order.

Polynomial partitioning and subsequent multiscale refinements brought
further progress on restriction \citep{Guth2016,WangBrooms2020}.
For the multiplier itself, Wu obtained the conjectured strict-order
range when $\max\{p,p'\}\ge13/4$ \citep{Wu2023}.
Guo, Oh, Wang, Wu, and Zhang developed a pseudoconformal approach
and an induction on two scales for transferring restriction methods
to Bochner--Riesz estimates \citep{GuoOhWangWuZhang2025}.
More recently, Gao, Wu, and Xi established the strict-order multiplier
range $\max\{p,p'\}\ge22/7$
\citep[Corollary~1.3]{GaoWuXi2026}. Wang and Wu's related diagonal
extension estimate holds for $p>22/7$
\citep[Theorem~0.2]{WangWu2024}. The connection also runs in the other
direction: Tao showed that the Bochner--Riesz conjecture implies the
spherical restriction conjecture \citep{Tao1999}. Our argument directly
establishes the scalar multiplier bound on $L^3(\RR^3)$ for every fixed
positive order.
Zipeng Wang has also proposed a proof of the Bochner--Riesz conjecture
in all dimensions \citep[version~6, Theorem~One]{ZipengWang2026}.

Let $S^2=\{\omega\in\RR^3:|\omega|=1\}$ and let $\sigma$ be surface
area measure on $S^2$. For bounded measurable $g:S^2\to\CC$, define
\[
 Eg(x)=\int_{S^2}g(\omega)e^{2\pi i x\cdot\omega}\,d\sigma(\omega).
\]
Tao's implication applied to \eqref{eq:intro-full-range} gives, for
every real $p>3$,
\[
 \norm{Eg}_{L^p(\RR^3)}
 \le C_p\norm{g}_{L^\infty(S^2,\sigma)},
\]
where $C_p<\infty$ depends only on $p$. This is the spherical extension
estimate also proved in \citet[Theorem~1.1]{OpenAISphere2026}.
For the separate diagonal extension theorem on compact smooth positively
curved surfaces, possibly with smooth boundary, see
\citet[Theorem~1.1]{OpenAIDiagonal2026}; that general-surface statement is
not deduced here from the spherical multiplier theorem.

An independent route through critical local smoothing gives the same
full strict fixed-radius range in \eqref{eq:intro-full-range}, uniformly
in the radius \citep[Corollary~11.2]{OpenAILocalSmoothing2026}.
It also gives strong maximal Bochner--Riesz bounds and almost-everywhere
summation for $3\le p<\infty$ and $\delta>1-3/p$
\citep[Theorem~11.1]{OpenAILocalSmoothing2026}.
These results are not inputs to the argument here.

The geometric input here is the sharp planar Furstenberg estimate of
Ren and Wang \citep{RenWang2025}. It belongs to a line of projection
theory that includes the classical work of Marstrand and Kaufman
\citep{Marstrand1954,Kaufman1968} and the discretized sum-product and
projection methods of Bourgain \citep{Bourgain2010}. Oberlin formulated
a sharp prediction for exceptional projection directions
\citep{Oberlin2012}. Important advances toward the Furstenberg problem
include the packing-dimension results of Orponen and Shmerkin
\citep{Orponen2020Packing,Shmerkin2022Packing}, their joint
Hausdorff-dimension improvement \citep{OrponenShmerkin2023Hausdorff},
and their proof of the conjectured bound when the associated family of
affine lines has equal Hausdorff and packing dimensions
\citep{OrponenShmerkin2026ABC}. Building on that regular case, Ren and
Wang removed the additional equality assumption. We use their finite
incidence theorem. Section~\ref{sec:projection} derives the form needed
for probability weights, auxiliary tags that may share a slope, and
incidence graphs of positive product mass; Section~\ref{sec:entropy}
supplies the finite conditioning used in later applications.

\subsection{From geometric growth to oscillatory cancellation}

The proof has two stages. We first control the information gained by
observing positive masses moving along lines. We then transfer that
control to wave packets, whose squared amplitudes are positive but
whose superposition is oscillatory. The main obstruction is at the
uncertainty boundary: an earlier packet cannot be assigned a position
as precisely as a later observation requires. The second stage records
this discrepancy as fine velocity information and shows that too much
growth would make that information exceed its finite capacity.

\paragraph{Positive masses on lines.}
A particle $T$ carries bounded data $(V,A)\in\RR^2\times\RR^2$ and
moves along $c\mapsto A+cV$. Fix a small parameter $\eps$.
An aperture observation at duration $t$ records velocity in a
rectangle of widths $\eps^b,\eps^a$, $b\le a$, and position in
a transported rectangle of widths $\eps^{b+t},\eps^{a+t}$.
Its orientation is also recorded. For $0<\gamma<1$, the weight
\[
 w(b,a)=a+b-\gamma(a-b)
\]
rewards fine resolution while penalizing eccentricity. The score in
\eqref{eq:classical-score} subtracts the particle information required
by the observation and adds twice the time entropy conditional on the
particle. Its largest asymptotic increase per unit duration is denoted
by $\betac$. Theorem~\ref{thm:classical-growth} proves
$\betac\le\gamma$.

To prove this bound, assume excessive growth and localize a sequence
whose scores approach the largest possible rate. The aspect penalty
makes the orientations of such observations coherent. Choosing the
least possible aspect produces a saturated path: its intermediate
observations attain the same growth rate. There are two alternatives.
An isotropic extremizer is excluded in Section~\ref{sec:isotropic} by
combining time-fiber richness, planar projections, and a
Loomis--Whitney argument. In the other alternative, the orientation
and longitudinal resolutions advance at a fixed ratio $m>0$.

The latter case produces four data coordinates with resolution speeds
$0,1,m,1+m$. Their particle-information profile is Lipschitz, but ordinary almost
everywhere differentiability does not give a derivative on the
particular plane selected by the path. Section~\ref{sec:characteristic-plane}
proves the needed trace and affine approximation there.
Planar projections then transfer lower bounds between the coordinate
rates. When $m=1$, the middle coordinates share a speed and only
their joint rate is available. Two conditional projection tests and a
reciprocal pinned-direction argument supply the missing inequalities.
The latter is related to the radial-projection bootstraps of Shmerkin
and Wang and of Orponen, Shmerkin, and Wang
\citep{ShmerkinWang2025Distance,OrponenShmerkinWang2024Radial}.
Section~\ref{sec:classical-closure} combines these rate constraints
with the saturated path's contact and jump conditions to obtain a
contradiction.

\paragraph{Oscillatory packets.}
Write $\lambda^{-1}=\eps^{\Planck}$. At duration $t$, a packet has
velocity depth $u_t=(\Planck-t)/2$ and position depth $u_t+t$.
The packet score in \eqref{eq:packet-score} uses squared oscillatory
amplitudes, and its largest growth rate is denoted by $\betaq$.
Sections~\ref{sec:packet-structure}--\ref{sec:packet-composition}
construct finite decompositions, control their tails, and track the
loss when only part of an assigned atom class is retained. These
constructions use the Fourier-series localization and almost
orthogonality underlying classical wave-packet methods
\citep{Tao2003,GuoOhWangWuZhang2025}.

Away from the uncertainty boundary, a coarser observation faithfully
locates the packets, so the positive-mass estimate applies. At the
boundary, a separate coherent estimate is required. It leaves a layer
of fine velocity information, which must be carried to the preceding
step. Section~\ref{sec:packet-repayment} proves that hypothetical
growth $\betaq>\gamma/2$ increases this information by at least
$(4\betaq-2\gamma)h$, up to controlled errors, at every backward
step of fixed length $2h$. The layer holds at most $2h$ units, up to vanishing errors. A fixed
finite number of steps therefore yields a contradiction and proves
Theorem~\ref{thm:packet-growth}: $\betaq\le\gamma/2$.

\paragraph{The multiplier.}
Packet growth gives the oscillatory estimate of
Theorem~\ref{thm:oscillatory-estimate}, with an arbitrarily small
positive power loss. On a physical shell of radius $\lambda$, the
Bochner--Riesz kernel has amplitude of order $\lambda^{-2-\delta}$.
The pseudoconformal distance-phase change of variables used in
\citet{GuoOhWangWuZhang2025} puts this shell into the required phase
class. Section~\ref{sec:multiplier} verifies all phase and amplitude
bounds uniformly in the shell and obtains the operator norm
$O_{\delta,\nu}(\lambda^{-\delta+\nu})$. Choosing $0<\nu<\delta$
proves Theorem~\ref{thm:bochner-riesz} by summation. A beta-integral
representation supplies the complex-order bounds used in Stein's
analytic interpolation theorem \citep{Stein1956}, completing
\eqref{eq:intro-full-range}.

\subsection{Finite laws and order of limits}

Entropy permits a common language for these geometric and analytic
steps. We use finite entropy and relative entropy
\citep{Shannon1948,KullbackLeibler1951}, conditional copies and
submodularity \citep{Tao2010Entropy}, and homogeneous profiles across
scales. Multiscale entropy and comparable-partition methods have
important precedents in \citet{HochmanShmerkin2012,Hochman2014}.
Section~\ref{sec:entropy} proves the precise finite-law statements
used here.

Small normalized mutual information yields a comparison with a product
law after trimming; it does not give exact independence.
Theorem~\ref{thm:finite-conditioning} specifies which labels are fixed
and which remain free, so later conditioning uses the correct comparison
population. For the characteristic-plane argument, first fix a
homogeneous component and its limiting profile, then choose the
approximation scale, and only afterwards take the finite precision
small relative to that scale.
Theorem~\ref{thm:characteristic-trace} and
Corollary~\ref{cor:finite-characteristic-trace} state these successive
steps.

Packet composition also keeps two finite populations distinct. Its
upper estimate concerns a retained superposition, while its earlier
lower test uses the whole assigned cap class.
Proposition~\ref{prop:finite-packet-composition} records the loss from
retention. Section~\ref{sec:packet-repayment} carries the remaining
fine velocity information through a fixed finite number of cuts,
using the one-step statement of Proposition~\ref{prop:finite-repayment}.

Throughout, \(\Ent,\Mut\) denote entropy and mutual information divided
by \(\log(1/\eps)\), whereas \(\Entn,\Mutn\) use natural logarithms.
The notation \(\logeps{r}\) uses the normalized units. The Planck depth
is always \(\Planck\), distinct from entropy. A positive depth increment
is fixed before the small-parameter limit; when increments subsequently
shrink, that additional limit is stated separately.
